\section{相关工作}
\label{sec:related-work-zh}

\subsection{异构图神经网络}

异构图神经网络将图表示学习扩展到包含多种节点类型和关系类型的图。HAN 使用节点级注意力和语义级注意力，从基于元路径的邻域中聚合信息。HGT 为异构图 Transformer 引入类型特定投影和注意力机制。MAGNN 同时建模元路径内聚合和元路径间聚合，以捕获细粒度语义依赖。HeCo 进一步利用网络模式视图和元路径视图之间的对比学习，提升异构图表示学习效果。

这些方法证明了语义聚合的价值，但它们通常面向干净图设计。在对抗性结构扰动下，它们对观测边和元路径诱导邻域的依赖可能导致污染语义被传播。DVCL 保留异构图学习的语义建模能力，同时显式引入拓扑净化和特征诱导视图以增强鲁棒性。

\subsection{鲁棒图神经网络}

鲁棒图神经网络旨在降低图模型对对抗性结构或噪声结构的敏感性。GCN-Jaccard 移除特征相似度较低节点之间的边。RGCN 使用高斯分布建模隐藏表示，以减弱对抗影响。Pro-GNN 结合稀疏和低秩约束联合学习图结构和 GNN 参数。GNNGuard 根据特征相似性为边分配权重，SimP-GCN 将相似性保持约束引入图卷积。

大多数鲁棒 GNN 是为同构图开发的。直接将它们应用到异构图可能忽略类型关系和语义路径。DVCL 的区别在于，它作用于元路径诱导的语义结构，并将特征诱导证据作为辅助视图，而不是仅作为简单的边剪枝规则。

\subsection{鲁棒异构图学习}

鲁棒异构图学习近年来受到关注，因为结构攻击可能通过异构语义被放大。RoHe、HeteGuard 和 HSeCo 等方法尝试通过关系感知扰动建模、过滤不可靠语义连接或改进语义拓扑学习来防御异构图神经网络。

这些方法与 DVCL 关系密切，但它们在构造候选邻域时仍然主要依赖受攻击的异构拓扑。DVCL 通过加入不直接依赖扰动边的特征诱导图来缓解这一限制。随后，跨视图目标鼓励拓扑视图表示和特征视图表示保持一致，从而在任一视图存在噪声时提升鲁棒性。

\subsection{图对比学习}

图对比学习通过最大化相关图视图之间的一致性来学习表示。DGI 将节点表示与图级摘要进行对比。MVGRL 对比扩散视图和邻接视图。GRACE 和 GCA 生成图增强，并在增强视图之间对齐节点表示。在异构图中，HeCo 对比网络模式视图和元路径视图以学习语义表示。

DVCL 同样基于跨视图一致性，但其视图设计目标是鲁棒性，而不仅是自监督表示学习。拓扑视图捕获净化后的异构语义，特征诱导视图提供抗攻击参考。因此，对比目标用于在结构攻击下提升语义拓扑和特征相似性之间的一致性。
